\noindent
La práctica calificada comprende la entrega (formato \verb|report|)
y la exposición (formato \verb|beamer|) de la resolución de los
ejercicios siguientes.
Fecha de entrega: lunes $13$ de abril de $2026$ a las 16 horas.
Referencia:~\cite{Atkinson2009}.

\section*{Aproximaciones de diferencias finitas}

\begin{questions}
	\question\label{question:1}

	Un enfoque para derivar fórmulas de diferencias para aproximar
	derivadas es el método de coeficientes indeterminados.
	Supongamos que $f$ es una función suave en $\mathbb{R}$.
	Sea $h>0$.
	Determine los coeficientes $a$, $b$ y $c$ de modo que
	\begin{equation*}
		af\left(x+h\right)+
		bf\left(x\right)+
		cf\left(x-h\right)
	\end{equation*}
	es una aproximación de $f^{\prime}\left(x\right)$ con un orden como
	lo más alto posible; es decir, elija $a$, $b$ y $c$ tales que
	\begin{equation*}
		\left|
		af\left(x+h\right)+bf\left(x\right)+
		cf\left(x-h\right)-f^{\prime}\left(x\right)
		\right|\leq
		\mathcal{O}\left(h^{p}\right).
	\end{equation*}

	\question

	Haz el mismo problema del Ejercicio~\ref{question:1} con
	$f^{\prime}\left(x\right)$ reemplazado por
	$f^{\prime\prime}\left(x\right)$.

	\question

	¿Es posible utilizar
	\begin{equation*}
		af\left(x+h\right)+
		bf\left(x\right)+
		cf\left(x-h\right)
	\end{equation*}
	con coeficientes elegidos adecuadamente para aproximar
	$f^{\prime\prime\prime}\left(x\right)$?
	¿Cuántos valores de función se necesitan para aproximar
	$f^{\prime\prime\prime}\left(x\right)$?

	\question\label{question:4}

	Para el problema del valor inicial de la ecuación de onda
	unidireccional
	\begin{equation}\label{eq:IVP-one-way-PDE}
		\begin{cases}\displaystyle
			\diffp{u}{t}=
			a\diffp{u}{x}+
			f   &
			\text{en }\mathbb{R}\times\left(0,\infty\right). \\
			u=f &
			\text{en }\mathbb{R}\times\left\{t=0\right\}.
		\end{cases}
	\end{equation}
	Donde $a\in\mathbb{R}$ es una constante, deriva alguna diferencia
	esquemas basados en varias combinaciones de aproximaciones en
	diferencias de la derivada temporal y la derivada espacial.

	\question

	La idea del esquema de Lax-Wendroff para resolver el problema del
	valor inicial~\eqref{eq:IVP-one-way-PDE} del
	Ejercicio~\ref{question:4} es el siguiente.
	Comience con la expansión de Taylor
	\begin{equation}\label{eq:taylor}
		u\left(x_{j},t_{m+1}\right)\approx
		u\left(x_{j},t_{m}\right)+
		h_{t}\diffp{u}{t}\left(x_{j},t_{m}\right)+
		\frac{h^{2}_{t}}{2}
		\diffp[2]{u}{t}\left(x_{j},t_{m}\right).
	\end{equation}
	De la ecuación diferencial tenemos
	\begin{equation*}
		\diffp{u}{t}=
		-a\diffp{u}{x}+
		f
	\end{equation*}
	y
	\begin{equation*}
		\diffp[2]{u}{t}=
		a^{2}\diffp[2]{u}{x}-
		a\diffp{f}{x}+
		\diffp{f}{t}.
	\end{equation*}
	Utilice estas relaciones para reemplazar las derivadas del tiempo
	en el lado derecho de~\eqref{eq:taylor}.
	Luego reemplace la primera y segunda derivadas espaciales por
	diferencias centrales.
	Finalmente, reemplace $\diffp{f}{x}$ por una diferencia central y
	$\diffp{f}{t}$ por una diferencia directa.
	Siga las instrucciones anteriores para derivar el esquema de
	Lax-Wendroff para resolver~\eqref{eq:IVP-one-way-PDE}.

	\question

	Dar esquemas hacia adelante, hacia atrás y Crank-Nicolson para el problema de valor inicial frontera
	\begin{equation*}
		\begin{cases}\displaystyle
			\diffp{u}{t}=
			\nu_{1}\diffp[2]{u}{x}+
			\nu_{2}\diffp[2]{u}{y}+
			f\left(x,y,t\right) &
			\text{en }\left(0,a\right)\times\left(0,b\right)\times\left(0,T\right). \\
			u=g_{l}             &
			\text{en }\left\{0\right\}\times\left[0,b\right]\times\left[0,T\right]. \\
			u=g_{r}             &
			\text{en }\left\{a\right\}\times\left[0,b\right]\times\left[0,T\right]. \\
			u=g_{b}             &
			\text{en }\left[0,a\right]\times\left\{0\right\}\times\left[0,T\right]. \\
			u=g_{t}             &
			\text{en }\left[0,a\right]\times\left\{b\right\}\times\left[0,T\right]. \\
			u=u_{0}             &
			\text{en }\left[0,a\right]\times\left[0,b\right]\times\left\{t=0\right\}.
		\end{cases}
	\end{equation*}
	Aquí, para los datos dados, $\nu_{1}$, $\nu_{2}$, $a$, $b$, $T$ son
	números positivos, y $f$, $g_{l}$, $g_{r}$, $g_{b}$, $g_{t}$,
	$u_{0}$ son funciones continuas de sus argumentos.

	\question

	Considere el problema del valor inicial frontera de una ecuación hiperbólica
	\begin{equation*}
		\begin{cases}\displaystyle
			\diffp[2]{u}{t}=
			\nu\diffp[2]{u}{x}+
			f\left(x,t\right)  &
			\text{en }\left(0,a\right)\times\left(0,T\right).   \\
			u=g_{1}            &
			\text{en }\left\{0\right\}\times\left[0,T\right].   \\
			u=g_{2}            &
			\text{en }\left\{a\right\}\times\left[0,T\right].   \\
			u=u_{0}            &
			\text{en }\left[0,a\right]\times\left\{t=0\right\}. \\
			\diffp{u}{t}=v_{0} &
			\text{en }\left[0,a\right]\times\left\{t=0\right\}.
		\end{cases}
	\end{equation*}
	Aquí, para los datos dados, $\nu$, $a$, $T$ son números positivos y
	$f$, $g_{1}$, $g_{2}$, $u_{0}$, $v_{0}$ son funciones continuas de sus argumentos.
	Analice cómo formar un esquema de diferencias finitas razonable para resolver el problema.
\end{questions}

\subsection*{Teorema de equivalencia de Lax-Richtmyer}

\begin{questions}
	\question

	Demuestre que el subespacio
	\vspace*{-.5\baselineskip}\setlength\belowdisplayshortskip{0pt}
	\begin{equation*}
		V_{0}\coloneqq
		\left\{
		v\in V \,\middle\vert\,
		v\left(x\right)=
		\sum_{j=1}^{n}a_{j}\sen\left(jx\right),
		a_{j}\in\mathbb{R},
		n\in\mathbb{N}
		\right\}
	\end{equation*}
	es denso en
	\begin{math}
		V\coloneqq
		\left\{
		v\in C\left(\left[0,\pi\right]\right)\mid
		v\left(0\right)=
		v\left(\pi\right)=0
		\right\}
	\end{math}.

	\question

	Analice el esquema de Crank-Nicolson para el problema~\eqref{eq:6.2.4}
	\begin{equation}\label{eq:6.2.4}
		\begin{cases}\displaystyle
			\diffp{u}{t}=
			\nu
			\diffp[2]{u}{x} &
			\text{en }\left[0,\pi\right]\times\left(0,T\right).   \\
			u=u_{0}         &
			\text{en }\left[0,\pi\right]\times\left\{t=0\right\}. \\
			u=0             &
			\text{en }\left\{0,\pi\right\}\times\left[0,T\right].
		\end{cases}
	\end{equation}

	\question

	Considere el esquema hacia adelante para resolver el problema de
	muestra~\eqref{eq:6.2.4}.
	Demuestre que $\left\|C\left(\Delta t\right)\right\|=\left|1-2r\right|+2r$,
	y $r\leq\frac{1}{2}$ es una condición necesaria y suficiente para
	estabilidad y convergencia.
\end{questions}

\subsection*{Más sobre convergencia y estabilidad}

\begin{questions}
	\question

	Sea $\left\{a,b,c\right\}\subset\mathbb{R}$ con $bc\geq 0$.
	Demuestre que los valores propios de la matriz
	\begin{equation*}
		Q=
		\begin{bmatrix}
			a      & c      & 0      & \cdots & 0 \\
			b      & a      & c      & \ddots & 0 \\
			0      & b      & \ddots & \ddots & 0 \\
			\vdots & \ddots & \ddots & a      & c \\
			0      & \cdots & 0      & b      & a
		\end{bmatrix}\in\mathbb{R}^{N\times N}
	\end{equation*}
	son
	\begin{equation*}
		\lambda_{j}=
		a+2\sqrt{bc}
		\cos\left(\frac{j\pi}{N+1}\right),
		\quad 1\leq j\leq N.
	\end{equation*}
	\textit{Sugerencia}:
	Para el caso no trivial $bc\neq 0$, escribe
	\begin{equation*}
		Q=
		aI+
		\sqrt{bc}D^{-1}\Lambda D
	\end{equation*}
	con $D$ es una matriz diagonal con elementos diagonales
	\begin{math}
		\sqrt{\frac{c}{b}},
		{\left(\sqrt{\frac{c}{b}}\right)}^{2},\dotsc,
		{\left(\sqrt{\frac{c}{b}}\right)}^{N}
	\end{math},
	y $\Lambda$ es la matriz tridiagonal con ceros en la diagonal y
	unos en la superdiagonal y subdiagonal,
	\begin{equation*}
		\Lambda=
		\begin{bmatrix}
			0      & 1      & 0      & \cdots & 0      \\
			1      & 0      & 1      & \ddots & \vdots \\
			0      & 1      & \ddots & \ddots & 0      \\
			\vdots & 0      & \ddots & 0      & 1      \\
			0      & \cdots & 0      & 1      & 0
		\end{bmatrix}\in\mathbb{R}^{N\times N}.
	\end{equation*}
	Luego encuentre los valores propios de $\Lambda$ siguiendo la definición
	del problema de valores propios y resolver un sistema de diferencias para
	componentes de vectores propios asociados.
	Un enfoque alternativo es relacionar la ecuación característica
	de $\Lambda$ a través de su fórmula de recursividad a polinomios de
	Chebyshev del segundo tipo~\cite[pág. 447]{Atkinson1991}.

	\question

	Proporcione un análisis de convergencia para el esquema de
	Crank-Nicolson para el problema~\eqref{eq:heatPDECrankNicolson}
	\begin{equation}\label{eq:heatPDECrankNicolson}
		\begin{cases}\displaystyle
			\frac{v^{m}_{j}-v^{m-1}_{j}}{h_{t}}=
			\nu
			\frac{
				\left(v^{m}_{j+1}-2v^{m}_{j}+v^{m}_{j-1}\right)+
				\left(v^{m-1}_{j+1}-2v^{m-1}_{j}+v^{m-1}_{j-1}\right)
			}{2h^{2}_{x}}+
			f^{m-\frac{1}{2}}_{j},             &
			1\leq j\leq N_{x}-1, 1\leq m\leq N_{t}. \\
			v^{m}_{0}=v^{m}_{N_{x}}=0,         &
			0\leq m\leq N_{t}.                      \\
			v^{0}_{j}=u_{0}\left(x_{j}\right), &
			0\leq j\leq N_{x}.
		\end{cases}
	\end{equation}
	en norma ${\left\|\cdot\right\|}_{2,h_{x}}$ aplicando el siguiente
	Teorema~\ref{thm:1},
	como se hace para los esquemas hacia adelante y hacia atrás en los ejemplos.
	\begin{theorem}\label{thm:1}
		Supongamos que el esquema~\eqref{eq:twolevels} es consistente y estable.
		\begin{equation}\label{eq:twolevels}
			\begin{aligned}
				\symbf{v}^{m+1} & =
				Q\symbf{v}^{m}+
				h_{t}\symbf{g}^{m},\quad 0\leq m\leq N_{t}-1, \\
				\symbf{v}^{0}   & =
				\symbf{u}_{0}.
			\end{aligned}
		\end{equation}
		Entonces el método es convergente.
		Además, si la solución $u$ es lo suficientemente suave como para que se cumpla
		\begin{equation}
			\sup_{m:mh_{t}\leq T}
			\left\|\symbf{\tau}^{m}\right\|\leq
			c\left(h^{p_{1}}_{x}+h^{p_{2}}_{t}\right),
		\end{equation}
		entonces tenemos la estimación del error
		\begin{equation*}
			\sup_{m:mh_{t}\leq T}
			\left\|\symbf{u}^{m}-\symbf{v}^{m}\right\|\leq
			c\left(h^{p_{1}}_{x}+h^{p_{2}}_{t}\right).
		\end{equation*}
	\end{theorem}

	\question

	Los esquemas hacia adelante, hacia atrás y Crank-Nicolson son
	casos particulares de una familia de esquemas de diferencias llamados
	métodos de punto medio generalizados.
	Sea $\theta\in\left[0,1\right]$ un parámetro.
	Entonces, para el valor inicial / frontera
	\begin{equation}\label{eq:heatPDE}
		\begin{cases}\displaystyle
			\diffp{u}{t}=
			\nu\diffp[2]{u}{x}+
			f\left(x,t\right) &
			\text{en }\left(0,\pi\right)\times\left(0,T\right).   \\
			u=0               &
			\text{en }\left\{0,\pi\right\}\times\left(0,T\right). \\
			u=u_{0}           &
			\text{en }\left(0,\pi\right)\times\left\{t=0\right\}.
		\end{cases}
	\end{equation}
	un esquema de punto medio generalizado es
	\begin{equation}\label{eq:heatPDETheta}
		\begin{cases}\displaystyle
			\begin{split}
				\frac{v^{m}_{j}-v^{m-1}_{j}}{h_{t}}=
				\nu\theta
				\frac{v^{m}_{j+1}-2v^{m}_{j}+v^{m}_{j-1}}{h^{2}_{x}}+
				\nu\left(1-\theta\right)
				\frac{v^{m-1}_{j+1}-2v^{m-1}_{j}+v^{m-1}_{j-1}}{h^{2}_{x}}+ \\
				\theta f^{m}_{j}+
				\left(1-\theta\right)f^{m-1}_{j},
			\end{split} &
			1\leq j\leq N_{x}-1, 1\leq m\leq N_{t}.                                                                                                                          \\
			v^{m}_{0}=v^{m}_{N_{x}}=0,                                                                                                                                     &
			0\leq m\leq N_{t}.                                                                                                                                               \\
			v^{0}_{j}=u_{0}\left(x_{j}\right),                                                                                                                             &
			0\leq j\leq N_{x}.
		\end{cases}
	\end{equation}
	Demuestre que para $\theta\in\left[\frac{1}{2},1\right]$,
	el esquema es incondicionalmente estable en las normas ${\left\|\cdot\right\|}_{2,h_{x}}$ y ${\left\|\cdot\right\|}_{\infty}$;
	para $\theta\in\left[0,\frac{1}{2}\right)$, el esquema es condicionalmente estable en la norma ${\left\|\cdot\right\|}_{2,h_{x}}$
	si $2\left(1-2\theta\right)r\leq 1$, y este es condicionalmente estable en la norma ${\left\|\cdot\right\|}_{\infty}$ si
	$2\left(1-\theta\right)r\leq 1$.
	Determinar los órdenes de convergencia de los esquemas.
\end{questions}

\section*{Sugerencia de lectura adicional}

Más detalles sobre el análisis teórico del método de diferencias finitas,
por ejemplo, tratamiento de otro tipo de condiciones de frontera,
dominios espaciales generales para problemas de dimensión espacial
superior, aproximación de problemas hiperbólicos o elípticos, se
pueden encontrar en varios libros sobre el tema, por ejemplo~\cite{Strikwerda2004}.
Para el método de diferencias finitas para problemas parabólicos,~\cite{Thomée1990}
es un análisis de encuesta en profundidad.
Otro enfoque popular para desarrollar métodos de diferencias finitas para problemas
parabólicos es el método de las líneas; vea~\cite[pág. 414]{Atkinson1991}
para una introducción que analiza algunos de los métodos de diferencias finitas de este
capítulo.
El libro~\cite{Gustafsson2008} proporciona una cobertura completa de
métodos de diferencias finitas de alto orden para EDPs dependientes del tiempo.

Para problemas de valor inicial / frontera de ecuaciones de evolución en alta dimensión espacial,
la estabilidad de un esquema explícito generalmente requiere que el tamaño del paso de tiempo
sea prohibitivamente pequeño.
Por otro lado, algunos esquemas implícitos son incondicionalmente
estables y el requisito de estabilidad no impone restricciones sobre el
tamaño del paso temporal.
La desventaja de un esquema implícito es que en cada nivel de tiempo
es posible que necesitemos resolver un sistema algebraico de muy gran escala.
La idea de la técnica de separación de operadores es dividir el
cálculo de cada paso de tiempo en varios subpasos de modo que cada
subpaso esté implícito solo en una variable espacial y al mismo
tiempo se mantenga una buena propiedad de estabilidad.
Los esquemas resultantes se denominan métodos de direcciones alternadas o métodos
de pasos fraccionario.
Véase~\cite{Marchuk1990,Yanenko1971} para discusiones detalladas.
Muchos fenómenos físicos se describen mediante leyes de conservación (conservación de masa, momento y energía).
Los métodos de diferencias finitas para las leyes de conservación constituyen un área de investigación grande.
El lector interesado puede vea~\cite{LeVeque1992} y~\cite{Godlewski2021}.

Los métodos de extrapolación son medios eficientes para acelerar la convergencia de soluciones numéricas.
Para métodos de extrapolación en el contexto del método de diferencias finitas, vea~\cite{Marchuk1983}.

\begin{questions}
	\question\label{question:research}

	Investigue sobre uno de los métodos de investigación, explique su idea principal, y discuta sus ventajas y limitaciones.
	Algunos ejemplos de métodos de investigación incluyen:
	\begin{itemize}
		\item Métodos de extrapolación para acelerar la convergencia de esquemas de diferencias finitas.
		\item Métodos de direcciones alternadas para resolver problemas de valor inicial / frontera de ecuaciones de evolución en alta dimensión espacial.
		\item Método multigrid para resolver sistemas lineales de gran escala que surgen de esquemas implícitos.
		\item Métodos de precondicionamiento iterativo para resolver sistemas lineales de gran escala que surgen de esquemas implícitos.
		\item Método de líneas para resolver problemas de valor inicial / frontera de ecuaciones de evolución.
		\item Método de separación de operadores para resolver problemas de valor inicial / frontera de ecuaciones de evolución.
		\item Método de paso fraccionario para resolver problemas de valor inicial / frontera de ecuaciones de evolución.
		\item Método de diferencia finita para la ecuación de Burgers.
	\end{itemize}
	% Considere el problema de valor inicial/frontera para una ecuación de evolución en
	% dimensión espacial superior $\mathbb{R}^{d}$ con $d\geq 2$:
	% \begin{equation*}
	% 	\begin{cases}\displaystyle
	% 		\diffp{u}{t}=
	% 		\sum_{i=1}^{d}\nu_{i}\diffp[2]{u}{x_{i}}+
	% 		f\left(\mathbf{x},t\right) &
	% 		\text{en }\Omega\times\left(0,T\right),\\
	% 		u=g &
	% 		\text{en }\partial\Omega\times\left[0,T\right],\\
	% 		u=u_{0} &
	% 		\text{en }\Omega\times\left\{t=0\right\},
	% 	\end{cases}
	% \end{equation*}
	% donde $\Omega\subset\mathbb{R}^{d}$ es un dominio espacial general, y $\nu_{i}>0$
	% son coeficientes de difusión.
	% Un esquema explícito para este problema requiere una condición de estabilidad
	% que impone una restricción severa sobre el tamaño del paso temporal $h_{t}$
	% (ver~\cite{Strikwerda2004,Gustafsson2008}).
	% En contraste, los esquemas implícitos (como Crank-Nicolson) son incondicionalmente
	% estables pero requieren resolver un sistema lineal de gran escala en cada paso temporal.

	% \textbf{Investigar lo siguiente:}
	% \begin{enumerate}
	% 	\item Compare el costo computacional de un esquema explícito versus un esquema
	% 	implícito para resolver este problema. Considere el número total de operaciones
	% 	aritméticas, el almacenamiento en memoria, y el tiempo de ejecución.
	% 	\item Explique cómo la técnica de separación de operadores (operador splitting)
	% 	o métodos de direcciones alternadas (ADI, alternating direction implicit)
	% 	mitiga este dilema. ¿Cómo se reduce la dimensionalidad efectiva del sistema
	% 	a resolver en cada subpaso?~\cite{Marchuk1990,Yanenko1971}.
	% 	\item Derive un esquema ADI para el problema anterior y analice su consistencia,
	% 	estabilidad y convergencia.
	% 	\item Discuta las limitaciones y ventajas del operador splitting respecto a otros
	% 	enfoques, como precondicionamiento iterativo de esquemas implícitos o uso de
	% 	multigrid.
	% \end{enumerate}

\end{questions}
